\section{学习通题目整理}
\subsection{第一章章节作业}
\begin{enumerate}
	\item 设 $A$、$B$、$C$ 为任意三个事件，用 $A$、$B$、$C$ 的运算表示下列事件：
	\begin{itemize}
		\item[(1)] 三个事件恰好发生两个；
		\item[(2)] 三个事件至少发生一个；
		\item[(3)] 三个事件最多发生两个；
		\item[(4)] $A$、$B$、$C$ 一个也不发生。
	\end{itemize}
	
	\vspace{5cm}
	
	\item 已知 $P(A) = 0.5$，$P(B) = 0.7$，$P(A \cup B) = 0.9$，求：
	\begin{itemize}
		\item[(1)] $P(A - B)$；
		\item[(2)] $P(\overline{A} \cup \overline{B})$。
	\end{itemize}
	
	\vspace{5cm}
	
	\item 将标号为 1, 2, 3, 4 的四个球随意地排成一行，求下列各事件的概率：
	\par (1) 各球自左至右或自右至左恰好排成 1, 2, 3, 4 的顺序；
	\par (2) 第 1 号球排在最右边或最左边；
	
	\vspace{5cm}
	
	\item 从 $[0, 1]$ 中随机地取两个数，求两数之积小于 $\dfrac{1}{4}$ 的概率。
	
	\vspace{5cm}
	
	\item 从 $1, 2, 3, \cdots, 15$ 这 15 个数中，甲乙两人各任取一个数（不重复），已知甲取到的数是 5 的倍数，求甲取到的数大于乙取到的数的概率。
	
	\vspace{5cm}
	
	\item 已知 $P(B \mid A) = P(A \mid B) = \dfrac{1}{2}$，$P(A) = \dfrac{1}{3}$，求 $P(A \cup B)$。
	
	\vspace{5cm}
	
	\item 根据以往的考试结果分析，努力学习的学生中可能有 90\% 考试及格，不努力学习的学生中可能有 90\% 考试不及格。据调查显示，学生中有 80\% 的人在努力学习，求：
	\par (1) 考试及格的学生是不努力学习的学生概率；
	\par (2) 考试不及格的学生是努力学习的学生概率。
	
	\vspace*{2cm}
	
	\item 设某人按如下原则决定某日的活动：如该天下雨则以 0.2 的概率外出购物，以 0.8 的概率去探访朋友；如该天不下雨，则以 0.9 的概率外出购物，以 0.1 的概率去探访朋友。又设某地下雨的概率是 0.3。试求：
	\par (1) 那天他外出购物的概率；
	\par (2) 若已知他那天外出购物，试求那天下雨的概率。
	
	\vspace{5cm}
	
	\item 设两两相互独立的三事件 $A, B, C$ 满足条件：$ABC = \varnothing$，$P(A) = P(B) = P(C) < 1/2$，且已知 $P(A \cup B \cup C) = 9/16$，求 $P(A)$。
	
	\vspace{5cm}
	
	\item 三架飞机（一架长机，两架僚机）一同飞往某目的地进行轰炸，但要达到目的地需要无线电导航，而只有长机有这种设备，到达目的地之前，必须经过敌方的高射炮阵地上空，这时任一飞机被击落的概率都是 0.2，到达目的地之后，各机将独立地进行轰炸，炸毁目标的概率都是 0.3，求目标被炸毁的概率。
\end{enumerate}

\newpage

\subsection{第二章章节作业}
\begin{enumerate}
	\item 做一系列独立试验，每次成功的概率为 $p(0 < p < 1)$，
	\begin{enumerate}
		\item 试验进行到成功与失败均出现时停止，求试验次数 $X$ 的分布律．
		\item 试验进行到出现 $n$ 次成功时停止，求已经失败的次数 $Y$ 的分布律．
	\end{enumerate}
	
	\vspace{5cm}
	
	\item 设离散型随机变量 $X$ 只取 $-1, 2, \pi$ 三个可能值，取各相应值的概率分别是 $a^2, -a$ 与 $a^2$，求 $X$ 的分布函数．
	
	\vspace{5cm}
	
	\item 已知离散型随机变量 $X$ 的分布函数为：
	\[
	F(x) = \begin{cases}
		0, & x < -1; \\
		0.3, & -1 \le x < 0; \\
		0.6, & 0 \le x < 1; \\
		0.8, & 1 \le x < 3; \\
		1, & x \ge 3.
	\end{cases}
	\]
	求 $X$ 的分布律，并计算 $P\{X < 1 \mid X \ne 0\}$．
	
	\vspace{5cm}
	
	\item 已知随机变量 $X \sim f(x) = a e^{-|x|}, -\infty < x < +\infty$，求
	\begin{enumerate}
		\item 参数 $a$；
		\item $X$ 的分布函数 $F(x)$．
	\end{enumerate}
	
	\vspace{5cm}
	
	\item 设连续型随机变量 $X$ 的分布函数为
	\[
	F(x) = \begin{cases}
		0, & x < 0 \\
		A x^2, & 0 \le x < 1 \\
		1, & x \ge 1
	\end{cases}
	\]
	试求 (1) 常数 $A$；(2) $X$ 落在 $\left(-1, \dfrac{3}{2}\right)$ 内的概率；(3) $X$ 的概率密度．
	
	\vspace{5cm}
	
	\item 已知随机变量 $X$ 的概率密度为
	\[
	f(x) = \begin{cases}
		ax + b, & 0 < x < 1, \\
		0, & \text{其他}.
	\end{cases}
	\]
	且 $P\left\{X > \dfrac{1}{2}\right\} = \dfrac{5}{8}$．(1) 求 $a, b$；(2) 计算 $P\left\{\dfrac{1}{4} < X \le \dfrac{1}{2}\right\}$．
	
	\vspace{5cm}
	
	\item 设随机变量 $X$ 在 $[2, 6]$ 上服从均匀分布，现对 $X$ 进行三次独立观察，试求至少有两次观测值大于 $3$ 的概率．
	
	\vspace{5cm}
	
	\item 设随机变量 $X$ 服从参数为 $\theta = 1$ 的指数分布，求关于 $x$ 的方程 $x^2 + Xx - X + 8 = 0$ 无实根的概率．
	
	\vspace{5cm}
	
	\item 设随机变量 $X$ 的分布律为
	\begin{center}
		\begin{tabular}{c|cccc}
			$X$ & $-2$ & $0$ & $2$ & $3$ \\ \hline
			$P$ & $\dfrac{1}{7}$ & $\dfrac{1}{7}$ & $\dfrac{3}{7}$ & $\dfrac{2}{7}$
		\end{tabular}
	\end{center}
	求 $Y = X^2$ 的分布律．
	
	\vspace{5cm}
	
	\item 已知随机变量 $X$ 的概率密度为
	\[
	f_X(x) = \begin{cases}
		e^{-x}, & x > 0, \\
		0, & x \le 0.
	\end{cases}
	\]
	求 $Y = e^X$ 的概率密度．
	
\end{enumerate}


























